% Rekurzivní výpočet 4! -- vlevo sestup (rozklad na menší podproblém),
% vpravo návrat (skládání výsledků).
\begin{tikzpicture}[x=1cm,y=1cm,font=\small]

    \node[task, minimum width=26mm] (f4) at (0, 0.00) {$4!=4\cdot 3!$};
    \node[task, minimum width=26mm] (f3) at (0,-1.40) {$3!=3\cdot 2!$};
    \node[task, minimum width=26mm] (f2) at (0,-2.80) {$2!=2\cdot 1!$};
    \node[task, minimum width=26mm] (f1) at (0,-4.20) {$1!=1$};

    % Sestup -- rozklad na menší podproblém
    \foreach \y in {-0.50,-1.90,-3.30}{
        \draw[diredge, darkred] (-0.7,\y) -- (-0.7,{\y-0.40});
    }
    \draw[diredge, darkred, line width=1.1pt] (-2.20,0.15) -- (-2.20,-4.35);
    \node[darkred, font=\scriptsize, rotate=90] at (-2.50,-2.10)
        {rozklad na menší podproblém};

    % Návrat -- skládání výsledků
    \foreach \y/\val in {-3.70/{$1$},-2.30/{$2$},-0.90/{$6$}}{
        \draw[diredge, darkgreen] (0.7,\y) -- (0.7,{\y+0.40})
            node[midway, right, darkgreen, font=\scriptsize] {\val};
    }
    \node[anchor=west, font=\small, darkgreen] at (1.20,0) {$=24$};
    \draw[diredge, darkgreen, line width=1.1pt] (2.20,-4.35) -- (2.20,0.15);
    \node[darkgreen, font=\scriptsize, rotate=270] at (2.50,-2.10)
        {skládání výsledků};

\end{tikzpicture}
